% ============================================================
% Section 1: Opening (Slides 1-3)
% ============================================================
\section{Opening}

% ----------------------------------------------------------
% Slide 1: Title Slide
% ----------------------------------------------------------
\begin{frame}[plain]
  \visualplaceholder[5cm]{Background gradient}

  \vspace{1.5cm}
  {\Huge\bfseries\textcolor{white}{The Story of Mathematics}}\\[0.6em]
  {\Large\textcolor{white!90}{5,000 Years of Human Genius}}\\[1.5em]
  {\large\textcolor{white!80}{From Clay Tablets to Artificial Intelligence}}\\[2em]
  {\normalsize\textcolor{white!70}{Prof.\ J.\ Osterrieder \hspace{2em} 2026}}

  \note{
    Welcome slide. Set the tone: this is a STORY, not a textbook chapter.
    We are going to meet the people behind the formulas.
    Mention the time span: roughly 5,000 years of documented mathematics
    (and even older if we count the Ishango bone at $\sim$20,000 BCE).
  }
\end{frame}

% ----------------------------------------------------------
% Slide 2: Why Study the History of Math?
% ----------------------------------------------------------
\begin{frame}[shrink=14]{Why Study the History of Mathematics?}

  \begin{columns}[T]
    \begin{column}{0.55\textwidth}
      \begin{itemize}
        \item Mathematics was \alert{invented by real people} solving real problems
        \item Every formula you have learned has a \textbf{human story} behind it
        \item Ideas emerged \textbf{independently} across continents
        \item Understanding \emph{where} math came from helps us understand \emph{what} it is
      \end{itemize}

      \vspace{0.8em}
      \begin{insightbox}
        Mathematics is not a fixed body of knowledge handed down from above
        --- it was created, debated, lost, and rediscovered by people just like you.
      \end{insightbox}
    \end{column}

    \begin{column}{0.42\textwidth}
      % World map with breakthrough locations
      \visualplaceholder[5cm]{Simplified world outline (rectangle placeholders for contine}
    \end{column}
  \end{columns}

  \note{
    Key message: math is a human endeavour, not a divine revelation.
    Point out the world map --- breakthroughs happened EVERYWHERE.
    ``Every formula you have learned has a human story behind it.''
    This lecture will take us from Central Africa 20,000 years ago
    to AI labs in 2026.
  }
\end{frame}

% ----------------------------------------------------------
% Slide 3: Our Journey Today -- Mega-Timeline
% ----------------------------------------------------------
\begin{frame}{Our Journey Today}

  \begin{center}
  \visualplaceholder[5cm]{Main timeline axis}
  \end{center}

  \vspace{0.3em}
  \begin{center}
    \small\textcolor{gray}{This timeline is our roadmap. We will mark ``\textcolor{accentOrange}{you are here}'' as we travel through each era.}
  \end{center}

  \note{
    Orientation slide. Walk students through the timeline left to right.
    Emphasise overlaps: Indian and Chinese math developed in PARALLEL with Greek math, not ``after'' it.
    Note the huge time span on the left (pre-classical) vs.\ the compressed modern era on the right.
    ``We will come back to this timeline at each section transition
    so you always know where we are in the story.''
  }
\end{frame}
